\chapter{KESIMPULAN DAN SARAN}
\label{bab:kesimpulan}

\section{Kesimpulan}

Dari tugas akhir ini dapat disimpulkan bahwa:
\begin{enumerate}[label=\arabic*.]
\item Program komputer optimasi beton bertulang pada struktur portal
      ruang yang dikembangkan dalam tugas akhir ini dengan menggunakan
      metoda optimasi polihedron fleksibel berhasil mengurangi harga
      struktur sekitar 9,9~\%.
\item Bertambahnya jumlah struktur desain sebagai titik coba
      menghasilkan struktur yang semakin murah, tetapi memerlukan waktu
      proses optimasi yang semakin lama.
\item Untuk jumlah variabel yang relatif besar, hasil yang didapatkan
      oleh metoda optimasi polihedron fleksibel kurang memuaskan, metoda
      polihedron fleksibel cukup baik hanya untuk kasus yang memiliki
      jumlah variabel desain yang relatif kecil.
\end{enumerate}

\section{Saran}

Untuk menyelesaikan masalah optimasi yang cukup rumit seperti optimasi
beton bertulang pada struktur portal ruang, diperlukan metoda optimasi
yang lebih baik lagi, ada beberapa saran yang mungkin dapat membantu
dalam mengembangkan optimasi beton bertulang pada struktur portal ruang
yaitu:
\begin{enumerate}[label=\arabic*.]
\item Dalam memilih bahasa pemrograman, sebaiknya digunakan bahasa yang
      dapat menembus memori lebih dari 64K, karena dalam pemrograman
      optimasi beton bertulang pada struktur portal ruang dibutuhkan
      dimensi sebesar mungkin, selain itu untuk sistem operasi dapat
      dicoba menggunakan sistem operasi yang sudah berbasiskan 64 bit
      atau yang \textit{flat memory}, misalnya menggunakan kernel Linux.
\item Metoda optimasi polihedron fleksibel perlu dikembangkan lagi,
      terutama dalam proses penelusurannya, sehingga metoda ini dapat
      menjadi lebih baik lagi.
\end{enumerate}
